\section{Typed SDKs and model-owned loops}\label{sec:C22:typed}
Typed SDKs place schemas and application dependencies around model interactions. Type validation can reject malformed arguments, constrain output shape, and expose errors before application code executes. Its guarantee is syntactic and representational. A valid customer identifier can still refer to the wrong customer, and a well-typed monetary amount can still exceed authorization. Semantic preconditions must therefore remain explicit application checks rather than be assumed consequences of schema adherence.

The inspected OpenAI Agents SDK separates an application-owned SDK path from managed runtime options~\cite{openai2026agents}. Handoffs, tool loops, session interfaces, and tracing provide useful abstractions, but applications still own tool effects and their storage semantics. Its SDK is provider-agnostic; the evaluation must identify the selected model adapter rather than infer it from the SDK name~\cite{executionproduct002}. PydanticAI makes typed interaction an inspectable design surface and also supplies a configurable harness, while Pydantic Graph and Pydantic Evals are separate packages~\cite{executionproduct004}. Restricting its inventory to output validation would omit its current context, tool and execution controls. A comparison should examine how validation failures are surfaced, whether repair attempts consume the same budget, and whether rejected calls enter the trajectory used for later training or evaluation.

The Claude Agent SDK occupies a different ownership boundary from a Messages API client. Current official documentation describes a Python/TypeScript library running the Claude Code binary in an application-operated process, with its agent loop, tool execution, permissions, sessions and hooks. A client SDK instead exposes model requests, while Claude Managed Agents supplies a hosted harness with a separately configured execution environment~\cite{cumb2026claudesdk,executionproduct061}. These are three distinct experimental treatments. A common model identifier does not make their context constructors or permission mechanisms equivalent.

One failure mode is repair-induced drift. A model proposes an invalid tool call, the SDK returns an error, and the next attempt changes a meaningful argument while correcting the format. If evaluation records only the final valid call, it omits the information that caused the transition. For a learned policy this can invalidate behavior-likelihood reconstruction. For a production service it can conceal repeated boundary probing. Validation and repair events should therefore be first-class trace records.

Model-owned loops choose their next action dynamically, whereas application-owned workflows constrain the available continuation structure. Neither is universally superior. Dynamic loops accommodate unforeseen subtasks but can cycle, invent unsupported tool names, or consume budgets without progress. Explicit workflows improve predictability but can fail when the task falls outside their branches. A hybrid design uses deterministic control for authority and commitment while permitting model-driven choice within a bounded action space.

The relevant baselines compare the same tool contract under direct function orchestration, a minimal agent loop, and the SDK's higher-level abstractions. Measure invalid calls, repair rounds, constraint violations, total cost, and completion. A framework's ergonomic advantage is real but requires a different study, such as developer implementation time under controlled tasks. It should not be inferred from task success or substituted for runtime reliability.

